This section describes the behaviour of squarefree monomials in the coinvariant ring, which is what the monomial computations of \Cref{sec: proofs of the main theorems} rest on. The essential feature is that the highest elementary symmetric function, $x_1 \cdots x_m = e_m$ vanishes in $R_m$: it is the only squarefree monomial that does, the annihilator of $x_A$ is generated by the complementary monomial $x_{A^c}$, and divisibility still detects membership in an ideal generated by one or two squarefree monomials.


\begin{lemma}
\label{lem:ann-squarefree-proofonly}
Suppose $x_A$ is a squarefree monomial in $\bK[x_1,\dots,x_m]$, where $A\subseteq [m]$. Then 
$\Ann_{R_m}(x_A)=(x_{A^c})$.
\end{lemma}

\begin{proof}
    Since $x_Ax_{A^c} = x_1\cdots x_m\in I_m$, then $(x_{A^c})\subseteq \Ann_{R_m}(x_A)$. 
    Conversely, given $f\in \Ann_{R_m}(x_A)$ it has a unique decomposition as a polynomial as $f = f_1 + f_2x_{A^c}$, where $f_1$ has no terms divisible by $x_{A^c}$. Then $x_Af_1\in I_m$ as well, and thus by hypothesis on $f_1$, we have $x_Af_1\in (e_1,\dots, e_{m-1})$. Indeed, otherwise if $x_Af_1$ has a nonzero term of the form $e_mg$ for some $g$, then $f_1$ has a nonzero term $x_{A^c}g$, contradicting our assumption on $f_1$.
    
    Since $e_1,\dots, e_{m-1},e_m$ is a regular sequence, $e_m$ is a nonzerodivisor on $S/(e_1,\dots, e_{m-1})$. Since $x_A|e_m$ as polynomials, then $x_A$ is also a nonzerodivisor on $S/(e_1,\dots, e_{m-1})$, thus `canceling' $x_A$ from $x_Af_1 = 0\mod (e_1,\dots, e_{m-1})$ gives $f_1\in (e_1,\dots, e_{m-1})$. Therefore, the image of $f_1$ in $R_m$ is zero. Hence, $f = f_2x_{A^c}\in (x_{A^c})$ as elements in $R_m$. 
\end{proof}

\begin{comment}
\begin{proof}
    Since $x_Ax_{A^c} = x_1\cdots x_m\in I_m$, then $(x_{A^c})\subseteq \Ann_{R_m}(x_A)$. 
    Conversely, let $f\in S$ represent an element of $\Ann_{R_m}(x_A),$ so that $x_Af\in I_m = (e_1,\ldots,e_{m-1})+(e_m).$
    Write $\widetilde{R}:= S/(e_1,\ldots,e_{m-1}).$
    Then in $\widetilde{R}$ we have that $x_Af = ce_m = cx_Ax_{A^c}$ for some $c\in\widetilde{r}.$
    Since $e_1,\dots, e_{m-1},e_m$ is a regular sequence, $e_m$ is a nonzerodivisor on $\widetilde{R}.$
    Since $x_A|e_m$ as polynomials, then $x_A$ is also a nonzerodivisor on $\widetilde{R}$, thus `canceling' $x_A$ from $x_A(f-cx_{A^c}) = 0$ in $\widetilde{R}$ gives $f = cx_{A^c}$ in $\widetilde{R}.$ 
    Since $R_m = \widetilde{R}/(e_m),$ the equality $f = cx_{A^c}$ also holds in $R_m,$ so $f\in(x_{A^c})$ as an element of $R_m.$
\end{proof}
\end{comment}

\begin{lemma}
\label{lem:quotient-and-annihilators}
\(R_m/(x_{i_1},\dots,x_{i_s})\) is the ring \(R_{m-s}\) up to relabeling the variables.
\end{lemma}

\begin{proof}
This follows by setting the chosen variables equal to zero in each
\(e_r\).  
\end{proof}

\begin{lemma}
\label{lem:contraction}
Let \(x_F,x_G\in \mathbb{C}[x_1,\ldots,x_m]\) be distinct squarefree monomials of degree at least two, with \(|F\cap G|\neq1\).  
For every squarefree \(x_H\in \mathbb{C}[x_1,\ldots,x_m]\),
\begin{enumerate}
    \item $ \overline{x_H}\in(\overline{x_F})\subseteq R_m
 \quad\Longleftrightarrow\quad
 x_F\mid x_H$
 \item $\overline{x_H}\in(\overline{x_F},\overline{x_G})  \quad\Longleftrightarrow\quad  x_F\mid x_H$ or $x_G\mid x_H.$
\end{enumerate}

\end{lemma}

\begin{proof}
For both statements, the reverse implication is trivially true.

To prove (1) by contrapositive, suppose \(x_F\nmid x_H\), i.e., inside $[m]$, $H^c \cap F\neq \varnothing$.


If \(|H^c \cap F|\geq 2\), choose \(a\in H^c \cap F\) and consider the quotient \(\pi_a: R_m \to R_m/(x_a)\).  
Then \(\pi_a(\overline{x_F})\) becomes zero, while \(\pi_a(\overline{x_H})\) remains a proper squarefree monomial because it still omits
the variables in \((H^c \cap F)\setminus\{a\}\). 
Thus \(\overline{x_H}\not\in(\overline{x_F})\subseteq R_m\).

If \(|H^c \cap F|=1\), i.e. $H^c\cap F = \{a\}$, write
\[ x_F=x_ax_I,\qquad x_H=x_Ix_J.\]
Here \(|I|\geq 1\), since \(\deg x_F\geq2\).
Without loss of generality, assume all $x_i$ appear in $H$ or $F$ so the complementary monomial to \(x_I\) is \(x_ax_J\). 
If \(\overline{x_H}\in (\overline{x_F})\), then there exists an $f$ such that \(\overline{x_H} = \overline{f}\overline{x_F}\), then
\[  \overline{x_I}\overline{(x_J-fx_a)}=\overline{0}. \]
\Cref{lem:ann-squarefree-proofonly} gives \(\overline{x_J-fx_a}\in(\overline{x_ax_J})\), and hence \(\overline{x_J}\in(\overline{x_a})\). 
Modding out by \(\overline{x_a}\) would make \(\overline{x_J}\) zero.  
This is impossible by \Cref{lem:ann-squarefree-proofonly} as \(\overline{x_J}\) is nonzero and proper there because it omits the nonempty factor \(\overline{x_I}\) and $x_a\not\vert x_J$ (as $x_F\not\vert x_H$).  
This proves Part (1).

Now suppose that neither \(x_F\) nor \(x_G\) divides \(x_H\).  
Choose an inclusion-minimal subset \(I\subseteq H^c\subset [m]\) meeting both \(F\) and \(G\).

If \(I\subsetneq H^c\), mod out $R_n$ by the variables indexed by \(I\).  
Both \(x_F\) and \(x_G\) vanish, while the image of \(x_H\) remains a proper squarefree monomial because it still omits every variable in \(H^c\setminus I\).  
Hence, \(\overline{x_H}\notin(\overline{x_F},\overline{x_G})\).


If \(I= H^c\), then either $I = \{a\} \subseteq F\cap G$ or $I = \{b,c\}$ where  $b\in F\setminus (F\cap G)$ and $c\in G\setminus (F\cap G)$.  
 
If $I= \{a\}$, $x_{F\cap G}=\gcd(x_F,x_G)$ has degree at least $2$ as $1\leq |F\cap G| \neq 1$.  
As $a\in F\cap G \cap I \subseteq H^c$, we have \(x_{F\cap G}\nmid x_H\).
Then Part 1 gives \(\overline{x_H}\notin (\overline{x_{F\cap G}})\supseteq(\overline{x_F},\overline{x_G})\).

If $I = \{b,c\}$, then mod out by $x_b$; label the quotient map $\pi_b$.
Then \(\pi_b(\overline{x_F})\) vanishes, while  \(\pi_b(\overline{x_G})\) remains squarefree of degree at least two and does not divide \(\pi_b(\overline{x_H})\), since \(x_c\) is missing.  
Part 1 shows that \(\pi_b(\overline{x_H})\not\in(\pi_b(\overline{x_G}))\).
Thus, \(\overline{x_H}\notin(\overline{x_F},\overline{x_G})\).

We finally see that in all cases,  \(\overline{x_H}\notin(\overline{x_F},\overline{x_G})\) so the contrapositive of the forward implication of Part (2) is proved.
\end{proof}


\begin{lemma}
\label{lem:endpoint-intersection-proofonly}
Let $A, B$ be disjoint subsets of $[m]$ such that $A\neq \emptyset$, and let $x_{A}$ and $x_{B}$ be the corresponding squarefree monomials in $\bK[x_1,\dots,x_m]$.
Then 
\begin{equation}\label{eq:ideal_int_xa}
    \bigcap_{a\in A}\left(x_{a},\;x_{A\setminus a}\,x_B\right) \subseteq (x_A,x_B)\qquad\text{inside }R_m.
\end{equation}
\end{lemma}

\begin{proof}
Induct on $|A|$.  
The case $|A|=1$ is immediate. 
For $|A|>1$, let $a_0$ be an arbitrary element of $A$ and write $x_A=x_{A\setminus a_0}\, x_{a_0}$. 

Given $f\in R_m$ that's in the intersection on the left-hand side of \eqref{eq:ideal_int_xa}, then in particular $f\in (x_{a_0},x_{A\setminus a_0}\,x_B)$, so
\[f=x_{a_0}g+x_{A\setminus a_0}\,x_Bh\]
for some $g,h$. 
For $a\in A\setminus a_0$, since $x_{a}$ divides the second term of this expansion, the condition $f\in(x_{a},x_{A\setminus a}\,x_B)$ gives $x_{a_0}g\in(x_{a},x_{A\setminus a}\,x_B)$. Reducing modulo $x_{a}$ gives
\[ 
\overline x_{a_0}\,\overline g\in (\overline{x_{A\cup B\setminus a}}) \subseteq R_{m-1}, 
\]
where $R_{m-1}$ is the coinvariant algebra with variables $x_i$ for $i\in [m]\setminus a$. Thus, there is some $\overline{\ell}\in R_{m-1}$ such that $\overline{x_{a_0}}\,\overline{g} = \overline{\ell}\,\overline{x_{A\cup B\setminus a}}$. Thus, $\overline{0} = \overline{x_{a_0}}(\overline{g}-\overline{\ell}\,\overline{x_{A\cup B\setminus\{a,a_0\}}})$ so $\overline{g}-\overline{\ell}\,\overline{x_{A\cup B\setminus\{a,a_0\}}}\in \Ann_{R_{m-1}}(\overline{x_{a_0}})$ and hence $\overline{g}-\overline{\ell}\,\overline{x_{A\cup B\setminus\{a,a_0\}}}\in (\overline{x_{[m]\setminus\{a,a_0\}}})$ by \Cref{lem:ann-squarefree-proofonly}. Therefore, $\overline{g}\in (\overline{x_{A\cup B\setminus \{a,a_0\}}})$. Lifting this last statement back to $R_m$, we have $g \in (x_a,x_{A\cup B\setminus \{a,a_0\}})$. Since this is for all $a\in A\setminus a_0$, we have $g\in \bigcap_{a\in A\setminus a_0} (x_a,x_{A\cup B\setminus \{a,a_0\}})$. Since $A\setminus a_0\neq \emptyset$, by our inductive hypothesis $g\in (x_{A\setminus a_0},x_B)$, which implies $f\in (x_A,x_B)$. Our induction is complete.
\end{proof}

\begin{lemma}
\label{lem:local-endpoint-colon-proofonly}
For $1\leq p<q\leq m+1$ with $p<m$,
\[(x_{[q,m]}:x_{[1,p]})\subseteq(x_{[p+1,q-1]},x_{[q,m]})\]
in $R_m$, where by convention $x_{\emptyset}=0$.
\end{lemma}

\begin{proof}
If $q=m+1$, 
\Cref{lem:ann-squarefree-proofonly} gives $(0:x_{[1,p]})=(x_{[p+1,m]})$ and we are done.

Now suppose $q\leq m$.
Let $f \in (x_{[q,m]}:x_{[1,p]})$, meaning that $fx_{[1,p]}\in (x_{[q,m]})$.
Quotienting by $x_a$ for some $a\geq q$ gives $\overline{f}\,\overline{x_{[1,p]}}\in (\overline{x_{[q,m]}}) =(\overline{0}) \subseteq R_{m-1}$, where $R_{m-1}$ is the coinvariant ring on variables $x_{a'}$ for $a'\in [m]\setminus a$. Thus, $\overline{f}\in \mathrm{Ann}_{R_{m-1}}(\overline{x_{[1,p]}})$
so by \Cref{lem:ann-squarefree-proofonly},
$\overline{f}\in(\overline{x_{[p+1,m]\setminus a}})$. Therefore, we have $f\in(x_a,x_{[p+1,m]\setminus a})$
for every $a\geq q$.  Applying \Cref{lem:endpoint-intersection-proofonly} with $A = [q,m]$ and $B = [p+1,q-1]$ then gives
$f\in(x_{[p+1,q-1]},x_{[q,m]})$.
\end{proof}

